% ---------------------------------------------------------------------------
% Chapitre 12 — Interpréter un CFG : worklist, arcs arrière, dominance
% Sources : notes/06-engine-fixpoint.md (§3.5, §3.8, §3.10, §3.11, §4.4,
% §4.9, §4.10, §6.4, §6.16, §8.9–§8.13, §8.17), src/engine/cfg_analyzer.rs,
% src/engine/dominance.rs, src/engine/eval.rs, src/engine/fixpoint.rs
% (appelants, exit_env, collect_thresholds), src/rules/api/query.rs
% (ExitDominance) ; ADR-014, ADR-025, ADR-042 §6.
% Sorties CLI observées au commit e67b10a avec target/debug/reactant sur
% /tmp/ch12/*.tsx. Déroulés pas à pas et ensembles de dominateurs relevés
% avec une sonde temporaire hors dépôt (/tmp/ch12/probe) : une réplique
% instrumentée de la boucle d'analyze_cfg, dont le résultat final a été
% comparé à celui d'analyze_cfg lui-même (identique), et des appels directs
% à compute_dominators, rpo et on_all_paths.
% ---------------------------------------------------------------------------
\chapter{Interpréter un CFG : worklist, arcs arrière, dominance}
\label{chap:cfg-analyzer-dominance}

\begin{objectifs}
  \item Savoir ce que calcule \code{analyze_cfg} sur un CFG : un environnement abstrait à la sortie de chaque bloc atteint et un store d'état enrichi de toutes les écritures de la passe. Savoir le dérouler à la main sur un bloc, un diamant et une boucle.
  \item Comprendre les trois opérateurs de fusion (première arrivée, join, \emph{widening} à seuils), le compteur d'arcs arrière tenu par en-tête de boucle, et pourquoi l'analyse d'un CFG termine.
  \item Connaître le raffinement de garde (\code{narrow_env_for_branch}) : les motifs reconnus, ce qu'ils ne reconnaissent pas, et pourquoi une garde réfutée n'élimine pas sa branche.
  \item Savoir relire une passe terminée (\code{exit_env}, \code{entry_env_of}) et sonder une expression après convergence (\code{eval_in_stores}, \code{ConvergedEval}).
  \item Connaître les questions \og sur tous les chemins\fg{} du moteur (\code{compute_dominators}, \code{DominatorTree}, \code{on_all_paths}, \code{ExitDominance}) et leurs conventions sur les sorties, les \code{throw} et les blocs inatteignables.
\end{objectifs}

\prerequis{\cref{chap:interpretation-abstraite} (\cref{def:treillis}, \cref{def:point-fixe}, \cref{def:widening}, \cref{def:must-may}, \cref{def:dominance}, \cref{alg:interpretation-abstraite-worklist}) ; \cref{chap:ir} (\cref{def:cfg}) ; \cref{chap:lowering-cfg} (\cref{def:lowering-back}, blocs de jonction) ; \cref{chap:treillis-simples} (intervalles, \code{widen_to}, raffinements) ; \cref{chap:transfert-stores} (\cref{def:environnement}, \cref{def:store}, \code{exec_stmt}).}

Le \cref{chap:transfert-stores} a donné la sémantique d'un \emph{pas} : que
vaut une expression, que fait une instruction. Le \cref{chap:point-fixe-composant}
donnera la boucle externe, celle qui enchaîne rendu, mémos, effets et handlers
jusqu'à ce que l'état React ne bouge plus. Entre les deux, il manque une pièce
de taille moyenne : exécuter abstraitement \emph{un corps entier}, c'est-à-dire
un CFG avec ses branchements et ses boucles. C'est le rôle de
\code{analyze_cfg}, dans \file{src/engine/cfg_analyzer.rs}. Le rendu d'un
composant, chacun de ses effets et chacun de ses handlers passent par elle, à
chaque tour de la boucle externe.

Le \cref{chap:interpretation-abstraite} en a donné l'idée
(\cref{alg:interpretation-abstraite-worklist}) et un premier déroulé. Ce
chapitre la reprend à la taille réelle du code, avec ses conventions : ce qui
est suivi bloc par bloc et ce qui ne l'est pas, où l'on joint, où l'on
accélère, ce que la condition d'un \code{if} apprend à l'analyse. Il présente
ensuite trois services que le moteur et les règles construisent sur le même
CFG une fois l'analyse finie : relire l'environnement d'un point du programme,
évaluer une expression contre les stores convergés, et répondre aux questions
de dominance (\og ce bloc s'exécute-t-il à chaque rendu ?\fg).

% ===========================================================================
\section{Le problème : un corps, plusieurs chemins}
\label{sec:cfg-analyzer-dominance-probleme}

Considérons un effet qui calcule une valeur par une boucle locale avant de
l'écrire dans l'état. C'est le composant \code{BoundedLocalLoop} de
\file{tests/fixtures/widening.tsx}, recopié sous un autre nom :

\begin{lstlisting}[language=TypeScript,caption={Un effet avec une boucle locale (\file{/tmp/ch12/loop.tsx})},label={lst:cfg-analyzer-dominance-loop}]
import { useState, useEffect } from "react";

export function LocalLoop() {
  const [total, setTotal] = useState(0);
  useEffect(() => {
    let i = 0;
    while (i < 5) {
      i = i + 1;
    }
    setTotal(i);
  }, []);
  return <div>{total}</div>;
}
\end{lstlisting}

Concrètement, l'effet s'exécute une fois (deps vides), la boucle fait cinq tours
et \code{setTotal} écrit 5. L'analyseur ne dit rien de ce composant :

\begin{sortie}
$ reactant check --verbose --info --no-color --project plain loop.tsx
[verbose] 1 components analyzed
[verbose] cache hits: 0, misses: 0
  [verbose] LocalLoop: 1 iteration(s), widened: []
   1 clean component(s) hidden, rerun with --show-clean
\end{sortie}

Pour arriver à ce silence, l'analyseur a dû répondre à une question qui n'a
rien de spécifique à React : \emph{quelles valeurs \code{i} peut-il avoir à
chaque point du corps de l'effet ?} Il ne peut pas exécuter la boucle (en
général, il ne connaît ni le nombre de tours ni les valeurs d'entrée). Il
calcule donc, pour chaque bloc du CFG, un environnement abstrait qui
sur-approxime toutes les valeurs concrètes possibles en ce point, sur tous les
chemins qui y mènent. Trois formes de graphe suffisent à tout expliquer :
\begin{itemize}
  \item la \emph{séquence} : un bloc passe son environnement de sortie au
    suivant, sans perte ;
  \item le \emph{diamant} : deux chemins se rejoignent, il faut un join ;
  \item la \emph{boucle} : un chemin revient sur lui-même, la suite des
    environnements peut croître sans fin, il faut un widening.
\end{itemize}
Les \cref{sec:cfg-analyzer-dominance-bloc,sec:cfg-analyzer-dominance-diamant,sec:cfg-analyzer-dominance-boucle}
les traitent dans cet ordre.

\subsection{Qui appelle \code{analyze_cfg}}

\code{analyze_cfg} n'est pas appelée par les règles : c'est une brique du
point fixe d'un composant. Ses cinq appelants sont tous dans
\code{analyze_component_impl} (\file{src/engine/fixpoint.rs}) :

\begin{table}[htbp]\centering\small
\begin{tabular}{@{}llll@{}}
\toprule
passe & CFG & environnement d'entrée & ligne \\
\midrule
rendu (chaque tour) & \code{render_cfg} & \code{initial_env} (props liées) & 366 \\
effets (chaque tour) & corps de chaque effet & \code{env_exit} du rendu & 429 \\
handlers (chaque tour) & corps de chaque handler & \code{env_exit} du rendu & 466 \\
rafraîchissement & \code{render_cfg} & \code{initial_env} & 548 \\
écritures pures & corps de chaque effet & \code{env_exit}, état $\Bot$ & 579 \\
\bottomrule
\end{tabular}
\caption{Les cinq appels d'\code{analyze_cfg} dans \file{src/engine/fixpoint.rs}.
Les deux derniers ont lieu après convergence ; ils sont décrits au
\cref{chap:point-fixe-composant}.}
\label{tab:cfg-analyzer-dominance-appelants}
\end{table}

Voici la signature :

\rustsnippet[label={lst:cfg-analyzer-dominance-sig}]{src/engine/cfg_analyzer.rs}{34}{46}{signature d'\code{analyze_cfg}}

Paramètre par paramètre :
\begin{itemize}
  \item \code{component} : l'identité du composant analysé. Elle ne sert
    qu'à construire l'\code{AnalysisCtx} de chaque bloc (la provenance des
    slots, \cref{chap:statevalue-stabilite}).
  \item \code{cfg} : le graphe à interpréter (\cref{def:cfg}).
  \item \code{entry_env} : l'environnement au début du bloc d'entrée.
  \item \code{state}, \code{memo} : les stores d'état et de mémos du tour
    courant de la boucle externe (\cref{def:store}), lus et non modifiés :
    la fonction en fait des copies.
  \item \code{transfer} : l'implémentation du trait \code{Transfer} ; en
    production c'est toujours \code{StateValueTransfer}.
  \item \code{widen_threshold}, \code{thresholds} : le nombre de passages par
    un arc arrière avant d'accélérer, et l'ensemble de seuils du widening
    (\cref{sec:cfg-analyzer-dominance-boucle}).
  \item \code{heap} : le tas abstrait, emprunté mutablement et partagé avec
    les autres passes du composant.
  \item \code{ctx}, \code{inter} : l'accès en lecture aux entrées de hooks
    (\code{QueryContext}) et le contexte inter-composants (\code{None} en
    analyse isolée).
\end{itemize}
Le résultat est un couple : une table bloc $\mapsto$ environnement, et le store
d'état enrichi des écritures de la passe.

% ===========================================================================
\section{Un bloc, puis une séquence}
\label{sec:cfg-analyzer-dominance-bloc}

\subsection{Deux tables, un store}

La fonction commence par préparer son état de travail :

\rustsnippet[label={lst:cfg-analyzer-dominance-init}]{src/engine/cfg_analyzer.rs}{47}{58}{initialisation}

\begin{definition}{Passe, environnement d'entrée, environnement de sortie}{cfg-analyzer-dominance-passe}
  Une \terme{passe} est une exécution d'\code{analyze_cfg} sur un CFG.
  Pendant une passe, chaque bloc $b$ atteint porte un \terme{environnement
  d'entrée} $\mathit{In}[b]$ (table \code{entry_envs}), valeur abstraite des
  variables avant sa première instruction, et un \terme{environnement de
  sortie} $\mathit{Out}[b]$ (table \code{exit_envs}), valeur après son
  dernier transfert. Seul $\mathit{Out}$ est rendu à l'appelant.
\end{definition}

Trois choix se lisent déjà dans ces douze lignes.
\begin{itemize}
  \item Seul le bloc d'entrée reçoit un environnement. Un bloc que rien
    n'atteint n'est jamais enfilé, donc jamais visité, et n'a pas de ligne
    dans \code{exit_envs} : l'absence de ligne joue le rôle de $\Bot$
    (\og aucun état concret n'arrive ici\fg).
  \item \code{state_out} est \emph{une seule} copie du store d'état pour toute
    la passe. L'état React n'est pas suivi bloc par bloc ; la
    \cref{sec:cfg-analyzer-dominance-etat} en tire les conséquences.
  \item La worklist est une file FIFO (\code{VecDeque}) dédupliquée par
    \code{in_worklist}. Ce n'est pas un parcours en ordre postfixe inverse :
    l'ordre de visite est l'ordre de découverte.
\end{itemize}

\begin{piege}
  Le type de retour s'appelle \code{BlockEnvs}, et sa documentation dit
  \og Per-block \emph{entry} environments\fg{}
  (\src{src/engine/cfg_analyzer.rs}{16}{17}). C'est faux : la fonction rend
  \code{exit_envs} (\srcline{src/engine/cfg_analyzer.rs}{123}), et sa propre
  documentation (\src{src/engine/cfg_analyzer.rs}{21}{22}) le dit
  correctement. Tout consommateur de \code{block_states} lit des
  environnements de \emph{sortie} ; celui qui veut l'entrée d'un bloc doit la
  reconstruire (\cref{sec:cfg-analyzer-dominance-relire}).
\end{piege}

\subsection{Le transfert d'un bloc}

La boucle principale défile un bloc, applique ses instructions et range le
résultat :

\rustsnippet[label={lst:cfg-analyzer-dominance-body}]{src/engine/cfg_analyzer.rs}{60}{89}{visite d'un bloc}

Étape par étape :
\begin{enumerate}
  \item \code{env_out} part d'une copie de $\mathit{In}[b]$.
  \item Le store de mémos est \emph{recopié à chaque bloc} (ligne 66) et la
    copie est jetée après le bloc : pendant une passe, les mémos sont figés.
    Ils ne sont recalculés qu'entre deux passes, par la boucle externe.
  \item Un \code{AnalysisCtx} est construit pour le bloc. Il emprunte
    mutablement \code{state_out} et le tas : ce sont les deux canaux par
    lesquels un bloc communique avec les autres \emph{sans} passer par les
    arêtes.
  \item Chaque instruction passe par \code{exec_stmt}
    (\cref{chap:transfert-stores}).
  \item Si le bloc se termine par \code{Return(e)}, les \emph{effets de bord}
    de \code{e} sont exécutés (paragraphe suivant).
  \item L'environnement obtenu devient $\mathit{Out}[b]$ ; une visite
    ultérieure du même bloc l'écrase.
\end{enumerate}

Les tests unitaires du fichier épinglent les cas les plus simples. Sur un CFG
d'un bloc vide, l'environnement d'entrée est rendu tel quel et l'état ne bouge
pas (\code{empty_cfg_entry_env_preserved},
\src{src/engine/cfg_analyzer.rs}{317}{336}). Un \code{let x = 42} donne
$\mathit{Out}[0](x) = [42,42]$ (\code{let_literal_extends_env},
\src{src/engine/cfg_analyzer.rs}{338}{363}). Une séquence de deux blocs,
\code{let x = {}} puis \code{return}, propage la valeur sans perte :
$\mathit{Out}[1](x)$ est une référence \code{PerRender}
(\code{two_block_exit_env_propagates},
\src{src/engine/cfg_analyzer.rs}{592}{649}). Sur un arc \code{Jump}, la
fonction de passage est l'identité.

\subsection{Les effets de bord d'un \code{Return}}
\label{sec:cfg-analyzer-dominance-return}

Un effet écrit avec une flèche concise n'a pas d'instruction :

\begin{lstlisting}[language=TypeScript,caption={Une flèche concise (\file{/tmp/ch12/concise.tsx})},label={lst:cfg-analyzer-dominance-concise}]
import { useState, useEffect } from "react";

export function Concise() {
  const [n, setN] = useState(0);
  useEffect(() => setN(n + 1));
  return <div>{n}</div>;
}
\end{lstlisting}

Le lowering abaisse le corps \code{() => setN(n + 1)} en un seul bloc sans
instruction, terminé par \code{Return(Call(setN, [n + 1]))}. Si
\code{analyze_cfg} n'exécutait que les instructions, l'appel du setter serait
invisible et cette boucle infinie passerait inaperçue : un faux négatif. D'où
les lignes 80 à 86 de \cref{lst:cfg-analyzer-dominance-body}, qui confient
l'expression du \code{Return} à \code{exec_expr_effects}, le même chemin que
celui d'une instruction-expression \code{e;}. Ce chemin
(\src{src/domains/interp/interpreter.rs}{114}{126}) fait trois choses : la
pré-passe des callbacks, l'appel de setter (mise à jour faible), et
l'évaluation d'une application de composant. La \emph{valeur} de retour n'est
pas calculée ici. Le diagnostic est bien produit :

\begin{sortie}
$ reactant check --verbose --no-color --project plain concise.tsx
  [verbose] Concise: 3 iteration(s), widened: [0]
  Concise  (2 hooks)  concise.tsx
    warn   infinite-loop  [hook:0]  (line 5:2)  this effect keeps pushing state `n` to new values on every run. Potential infinite render loop
\end{sortie}

Trois tests de \file{src/engine/cfg_analyzer.rs} fixent le comportement :
\begin{itemize}
  \item un setter dans un \code{Return} écrit (l.~850--893,\\
    \code{setter_in_return_terminator_updates_state}) ;
  \item un \code{Return(Lit(Unit))} ne fait rien (l.~896--937,
    \code{unit_return_terminator_is_noop}) ;
  \item un \emph{functional updater} passé dans un \code{Return} est exécuté
    (l.~940--1018,\\ \code{functional_updater_in_return_terminator_fires}).
\end{itemize}
Le deuxième cas est le plus fréquent : depuis l'\adr{25}, la fin d'un corps
sans \code{return} est scellée \code{Return(Lit(Unit))}.

\begin{remarque}
  Le commentaire des lignes 80 à 83 rappelle une correction : l'ancien code
  fabriquait une instruction \code{Stmt::ExprStmt} jetable pour réutiliser le
  chemin des instructions. Le trait \code{Transfer} expose désormais
  \code{exec_expr_effects} comme la définition unique de \og exécuter une
  expression pour ses effets\fg{} (\src{src/domains/mod.rs}{133}{144}).
\end{remarque}

\subsection{L'environnement de fin d'un corps}

Les passes d'effets et de handlers partent de \code{env_exit}, l'environnement
de fin du rendu. Il est calculé à partir des sorties de la passe de rendu :

\rustsnippet[label={lst:cfg-analyzer-dominance-exit-env}]{src/engine/fixpoint.rs}{1226}{1237}{l'environnement de sortie d'un CFG}

C'est le join des $\mathit{Out}[b]$ des blocs terminés par un \code{Return}.
Un \code{Return} inatteignable n'a pas de ligne dans \code{block_states} et le
\code{filter_map} l'écarte. La réduction se fait par \code{reduce}, pas par un
\code{fold} partant de \code{AbstractEnv::bottom()}, et c'est voulu : le
paragraphe suivant montre que la table vide n'est pas neutre pour le join. La
même fonction existe en méthode publique sur \code{AnalysisResult}
(\src{src/engine/analysis_result.rs}{279}{288}) pour les règles.

% ===========================================================================
\section{Le diamant : joindre aux points de jonction}
\label{sec:cfg-analyzer-dominance-diamant}

\subsection{Ce que reçoit chaque successeur}

Après la visite d'un bloc, \code{analyze_cfg} demande à \code{outgoing} la
liste de ses successeurs, chacun avec l'environnement qu'il reçoit :

\rustsnippet[label={lst:cfg-analyzer-dominance-outgoing}]{src/engine/cfg_analyzer.rs}{126}{150}{l'environnement passé à chaque successeur}

\code{outgoing} lit le \emph{terminateur} du bloc, pas la liste
\code{cfg.edges}. Un \code{Jump} passe l'environnement tel quel ; un
\code{Branch} le raffine selon sa condition, vers \code{then_} avec
\code{taken = true} et vers \code{else_} avec \code{taken = false}
(\cref{sec:cfg-analyzer-dominance-gardes}) ; \code{Return} et
\code{Unreachable} n'ont pas de successeur. La branche finale, qui lit
\code{cfg.successors}, ne sert que pour un identifiant de bloc absent de
\code{cfg.blocks}. Sur un CFG bien formé, les deux lectures coïncident :
\code{CFG::validate} exige que toute cible de terminateur ait son arête
(\src{src/ir/cfg.rs}{181}{227}).

\subsection{Un diamant, pas à pas}

Le test \code{diamond_joins_at_merge_point}
(\src{src/engine/cfg_analyzer.rs}{651}{753}) construit à la main l'IR de ce
fragment :

\begin{lstlisting}[language=TypeScript,caption={Le diamant du test, en source équivalent},label={lst:cfg-analyzer-dominance-diamant}]
let x;
if (true) { x = 1; } else { x = {}; }
// ici : x vaut 1 ou un objet frais
\end{lstlisting}

\begin{figure}[htbp]\centering
\begin{tikzpicture}[node distance=9mm and 14mm]
  \node[cfgbloc] (b0) {B0 (entrée)\\branch true};
  \node[cfgbloc,below left=of b0] (b1) {B1\\let x = 1\\jump B3};
  \node[cfgbloc,below right=of b0] (b2) {B2\\let x = \{\}\\jump B3};
  \node[cfgbloc,below=28mm of b0] (b3) {B3\\return};
  \draw[fleche] (b0) -- node[etiquette,left]{IfTrue} (b1);
  \draw[fleche] (b0) -- node[etiquette,right]{IfFalse} (b2);
  \draw[fleche] (b1) -- (b3);
  \draw[fleche] (b2) -- (b3);
  \node[etiquette,left=2mm of b1,align=right] {$\mathit{Out}$ : $x \mapsto$ num $[1,1]$};
  \node[etiquette,right=2mm of b2,align=left] {$\mathit{Out}$ : $x \mapsto$ ref \code{PerRender}};
  \node[etiquette,below=2mm of b3,align=center] {$\mathit{In} = \mathit{Out}$ : $x \mapsto$ num $[1,1]$ \emph{et} ref \code{PerRender}\\(join du produit, pas $\Top$)};
  \node[etiquette,right=2mm of b0,align=left] {$\mathit{Out}$ : $\{\}$ (table vide)};
\end{tikzpicture}
\caption{Le diamant de \code{diamond_joins_at_merge_point}, annoté des
environnements calculés par \code{analyze_cfg}.}
\label{fig:cfg-analyzer-dominance-diamant}
\end{figure}

La \cref{fig:cfg-analyzer-dominance-diamant} montre le déroulé. La condition
est le littéral \code{true} : \code{narrow_env_for_branch} ne reconnaît pas ce
motif (\cref{sec:cfg-analyzer-dominance-gardes}), les deux branches reçoivent
la table vide inchangée et sont visitées. La branche \code{else} est pourtant
concrètement morte ; l'analyse l'explore quand même, ce qui est sound. B1 est
visité en premier et installe $\mathit{In}[B3]$ par \og première
arrivée\fg{} (\code{None => outgoing_env}, ligne 98). B2 arrive ensuite et
provoque le join (ligne 109). Le test vérifie que le join du produit
\code{StateValue} garde \emph{les deux} sortes : la composante numérique vaut
$[1,1]$, la composante référence \code{PerRender}, et la valeur n'est pas
$\Top$ (\cref{def:statevalue}). Un domaine plat aurait perdu toute information
ici, c'était la motivation de l'ADR-015.

\subsection{Le join des environnements, et pourquoi $\Bot$ n'y est pas neutre}

Le join de deux \code{AbstractEnv} est point à point, avec une règle
particulière pour une variable présente d'un seul côté :

\rustsnippet[label={lst:cfg-analyzer-dominance-env-join}]{src/domains/stores/abstract_env.rs}{130}{142}{join des valeurs : une variable d'un seul côté vaut $\Top$}

La règle suit l'\cref{inv:env-top} : une variable absente d'un environnement
s'y lit $\Top$, donc son join avec n'importe quoi vaut $\Top$. Dans le diamant
ci-dessus, si seule la branche \code{then} liait \code{x}, B3 lirait
\code{x} à $\Top$. C'est sound : sur le chemin \code{else}, \code{x} a la
valeur qu'il avait avant le \code{if}, que la table ne connaît pas.

\begin{piege}
  La conséquence est que \code{AbstractEnv::bottom()}, la table vide, n'est
  \emph{pas} un élément neutre du join : $\Bot \join \rho$ rend toutes les
  variables de $\rho$ à $\Top$. Écrire \code{fold(AbstractEnv::bottom(), join)}
  sur une liste d'environnements détruit toute l'information ; c'est pourquoi
  \code{exit_env} utilise \code{reduce}. De même, la première arrivée sur un
  bloc \emph{installe} l'environnement reçu au lieu de le joindre à une
  entrée vide. En revanche l'ordre \code{leq} de la table lit une variable
  absente comme $\Bot$ (\src{src/domains/stores/abstract_env.rs}{216}{239}) :
  les deux conventions ne sont pas duales, et \code{analyze_cfg} ne se sert
  d'ailleurs pas de \code{leq}.
\end{piege}

\subsection{Détecter un changement}

Une fois la nouvelle entrée calculée, le successeur n'est ré-enfilé que si elle
diffère de l'ancienne (ligne 114, \cref{lst:cfg-analyzer-dominance-propag}
plus bas). Le test est une égalité structurelle \code{!=}, pas
$\neg(\mathit{new} \lleq \mathit{old})$. Les deux sont équivalents ici, parce
que la nouvelle entrée est toujours au-dessus de l'ancienne (un join ou un
widening de l'ancienne) : pour $\mathit{new} \lgeq \mathit{old}$, on a
$\mathit{new} \lleq \mathit{old} \iff \mathit{new} = \mathit{old}$ dans un
ordre partiel.

\begin{remarque}
  L'égalité structurelle voit parfois un changement là où il n'y en a pas.
  Un intervalle vide n'a pas de représentation canonique : \code{Interval}
  compare ses bornes (\src{src/domains/impls/interval.rs}{23}{27}), et
  \code{narrow_geq(5)} appliqué à $[0,0]$ rend l'intervalle $\{lo = 5, hi =
  0\}$, vide, tandis que le même raffinement de $[0,1]$ rend $\{lo = 5, hi =
  1\}$, vide aussi, mais différent pour \code{==}. La sonde du chapitre
  l'observe sur la boucle de la section suivante : la sortie de boucle, dont
  l'entrée reste $\Bot$, est ré-enfilée et revisitée à chaque tour de la
  boucle (pas 7 et 10 du \cref{tab:cfg-analyzer-dominance-iterations}). C'est
  un surcoût, borné par le nombre de changements de l'en-tête, sans effet
  sur le résultat.
\end{remarque}

% ===========================================================================
\section{La boucle : arcs arrière, compteur par en-tête, widening à seuils}
\label{sec:cfg-analyzer-dominance-boucle}

\subsection{La propagation}

Voici la fin de la boucle principale, qui traite chaque successeur :

\rustsnippet[label={lst:cfg-analyzer-dominance-propag}]{src/engine/cfg_analyzer.rs}{91}{121}{fusion et propagation vers les successeurs}

Trois cas pour la nouvelle entrée $\nu$ d'un successeur $s$ recevant $\rho$ :
\begin{itemize}
  \item $s$ n'a pas encore d'entrée : $\nu = \rho$ ;
  \item l'arc $b \to s$ est de nature \code{EdgeKind::Back} : on incrémente le
    compteur de $s$ ; tant qu'il est sous \code{widen_threshold},
    $\nu = \mathit{In}[s] \join \rho$, ensuite
    $\nu = \mathit{In}[s] \widen_T \rho$ (\code{widen_to} avec les seuils $T$) ;
  \item sinon : $\nu = \mathit{In}[s] \join \rho$.
\end{itemize}

Le compteur \code{back_edge_counts} est indexé par le bloc \emph{cible},
c'est-à-dire l'en-tête de la boucle, et non par l'arc. Deux arcs arrière vers
le même en-tête (la fin du corps et un \code{continue}) incrémentent le même
compteur : une boucle qui contient un \code{continue} peut atteindre le seuil
en moins de tours. Le compteur compte les \emph{arrivées}, qu'elles changent
l'entrée ou non. Il est local à la passe : chaque appel à \code{analyze_cfg}
repart de zéro. Enfin, sa valeur de déclenchement \code{widen_threshold} est
le même paramètre que celui de la boucle externe, 3 par défaut
(\src{src/engine/fixpoint.rs}{51}{60}).

La nature de l'arc est retrouvée en parcourant \code{cfg.edges} (lignes 92 à
95) : c'est le lowering qui décide où l'on accélère, en marquant \code{Back}
les arêtes qui reviennent vers un en-tête (\cref{def:lowering-back}).

\subsection{L'ensemble des seuils}

Les seuils viennent du programme lui-même. Ils sont récoltés une fois par
composant, avant la boucle externe :

\rustsnippet[label={lst:cfg-analyzer-dominance-thresholds}]{src/engine/fixpoint.rs}{1185}{1206}{récolte des seuils du widening}

Tous les littéraux numériques du rendu, des corps d'effets et de handlers
(y compris ceux des fonctions imbriquées, \src{src/engine/fixpoint.rs}{1212}{1224})
et des initialiseurs de \code{useState} entrent dans $T$, triés et dédupliqués.
Le commentaire donne la justification : récolter trop est sans danger,
l'ensemble reste fini et un seuil de plus n'ajoute qu'une borne candidate. Le
même $T$ est passé à toutes les passes du composant, donc aussi aux boucles
locales. L'opérateur \code{widen_to} lui-même est présenté au
\cref{chap:treillis-simples} (\cref{thm:treillis-simples-widen-to}) : une
borne qui croît saute au plus proche seuil qui l'englobe, ou à $\pm\infty$
s'il n'y en a pas. Sur \code{StateValue}, seul le composant numérique utilise
les seuils ; les autres composantes prennent le widening ordinaire
(\src{src/domains/impls/state_value.rs}{670}{676}).

\subsection{Le compteur de boucle, pas à pas}

Deux tests de \file{src/engine/cfg_analyzer.rs} (l.~450--509),
\begin{itemize}
  \item \code{loop_counter_unbounded_without_thresholds},
  \item \code{loop_counter_bounded_by_threshold},
\end{itemize}
analysent le même CFG, \code{counting_loop_cfg} (l.~365--448), l'IR de
\code{let i = 0; while (i < 5) { i = i + 1; }}. Le premier passe
$T = \emptyset$, le second $T = \{5\}$ (la constante de la garde, comme la
récolterait \code{collect_thresholds}).

\begin{figure}[htbp]\centering
\begin{tikzpicture}[node distance=10mm and 16mm]
  \node[cfgbloc] (b0) {B0 (entrée)\\let i = 0\\jump B1};
  \node[cfgbloc,below=of b0] (b1) {B1 (en-tête)\\branch i < 5};
  \node[cfgbloc,below left=of b1] (b2) {B2 (corps)\\i = i + 1\\jump B1};
  \node[cfgbloc,below right=of b1] (b3) {B3 (sortie)\\return undefined};
  \draw[fleche] (b0) -- node[etiquette,right]{Unconditional} (b1);
  \draw[fleche] (b1) -- node[etiquette,above left]{IfTrue} (b2);
  \draw[fleche] (b1) -- node[etiquette,above right]{IfFalse} (b3);
  \draw[flechemay] (b2.west) to[bend left=50] node[etiquette,left]{Back} (b1.west);
  \node[etiquette,right=2mm of b0,align=left] {$\mathit{Out}$ : $i = [0,0]$};
  \node[etiquette,right=2mm of b1,align=left] {$\mathit{Out}$ : $[0,+\infty]$ ($T=\emptyset$)\\\phantom{$\mathit{Out}$ : }$[0,5]$ ($T=\{5\}$)};
  \node[etiquette,below=2mm of b2,align=center] {$\mathit{In}$ : $[0,4]$, $\mathit{Out}$ : $[1,5]$};
  \node[etiquette,below=2mm of b3,align=center] {$\mathit{Out}$ : $[5,+\infty]$ ($T=\emptyset$)\\$[5,5]$ ($T=\{5\}$)};
\end{tikzpicture}
\caption{\code{counting_loop_cfg} et ses environnements finaux (composante
numérique de \code{i}). L'arc arrière est en pointillé.}
\label{fig:cfg-analyzer-dominance-boucle}
\end{figure}

Le \cref{tab:cfg-analyzer-dominance-iterations} donne chaque visite. Il a été
relevé avec la réplique instrumentée de la sonde ; les environnements finaux
qu'elle trouve sont identiques à ceux d'\code{analyze_cfg} et à ceux que les
deux tests vérifient.

\begin{table}[htbp]\centering\small
\begin{tabular}{@{}rlll@{}}
\toprule
pas & bloc & $\mathit{In} \to \mathit{Out}$ (valeur de \code{i}) & mises à jour des successeurs \\
\midrule
\multicolumn{4}{@{}l}{\emph{pas communs aux deux variantes}} \\
1 & B0 & $\Top \to [0,0]$ & $\mathit{In}[B1] := [0,0]$ (première arrivée) \\
2 & B1 & $[0,0]$ & $\mathit{In}[B2] := [0,0]$ ; $\mathit{In}[B3] := \Bot$ \\
3 & B2 & $[0,0] \to [1,1]$ & arrière n\textsuperscript{o}\,1, join : $\mathit{In}[B1] := [0,1]$ \\
4 & B3 & $\Bot$ & -- \\
5 & B1 & $[0,1]$ & $\mathit{In}[B2] := [0,1]$ ; $\mathit{In}[B3]$ reste vide (ré-enfilé) \\
6 & B2 & $[0,1] \to [1,2]$ & arrière n\textsuperscript{o}\,2, join : $\mathit{In}[B1] := [0,2]$ \\
7 & B3 & $\Bot$ & -- \\
8 & B1 & $[0,2]$ & $\mathit{In}[B2] := [0,2]$ ; $\mathit{In}[B3]$ reste vide (ré-enfilé) \\
\midrule
\multicolumn{4}{@{}l}{\emph{sans seuil, $T = \emptyset$}} \\
9 & B2 & $[0,2] \to [1,3]$ & arrière n\textsuperscript{o}\,3, widening : $[0,2] \widen [1,3] = [0,+\infty]$ \\
10 & B3 & $\Bot$ & -- \\
11 & B1 & $[0,+\infty]$ & $\mathit{In}[B2] := [0,4]$ ; $\mathit{In}[B3] := [5,+\infty]$ \\
12 & B2 & $[0,4] \to [1,5]$ & arrière n\textsuperscript{o}\,4 : $[0,+\infty] \widen [1,5]$ inchangé \\
13 & B3 & $[5,+\infty]$ & -- (file vide) \\
\midrule
\multicolumn{4}{@{}l}{\emph{avec seuil, $T = \{5\}$}} \\
9 & B2 & $[0,2] \to [1,3]$ & arrière n\textsuperscript{o}\,3, widening : $[0,2] \widen_T [1,3] = [0,5]$ \\
10 & B3 & $\Bot$ & -- \\
11 & B1 & $[0,5]$ & $\mathit{In}[B2] := [0,4]$ ; $\mathit{In}[B3] := [5,5]$ \\
12 & B2 & $[0,4] \to [1,5]$ & arrière n\textsuperscript{o}\,4 : $[0,5] \widen_T [1,5]$ inchangé \\
13 & B3 & $[5,5]$ & -- (file vide) \\
\bottomrule
\end{tabular}
\caption{Les treize visites d'\code{analyze_cfg} sur \code{counting_loop_cfg}
(\code{widen_threshold = 3}). Seule la composante numérique de \code{i} est
montrée ; $\Bot$ désigne un intervalle vide.}
\label{tab:cfg-analyzer-dominance-iterations}
\end{table}

Quelques points méritent d'être soulignés.
\begin{itemize}
  \item Au pas 1, \code{i} n'est pas lié dans la table d'entrée : sa lecture
    vaut $\Top$ (\cref{inv:env-top}). Le \code{let} l'écrase.
  \item Au pas 2, la branche \og faux\fg{} raffine $[0,0]$ par
    \code{i >= 5} et obtient un intervalle vide. Le bloc B3 est quand même
    enfilé et visité (pas 4) : un raffinement vers $\Bot$ ne rend pas un
    bloc inatteignable (\cref{sec:cfg-analyzer-dominance-gardes}). Ici c'est
    sans conséquence, B3 ne fait rien.
  \item Les deux premiers retours par l'arc arrière joignent : l'en-tête voit
    $[0,1]$ puis $[0,2]$, exactement les valeurs concrètes des deux premiers
    tours. Le troisième retour déclenche le widening.
  \item Sans seuil, la borne supérieure saute à $+\infty$. Le raffinement de
    la garde rattrape une partie de la perte \emph{sur les arcs sortants} :
    le corps reçoit $[0,4]$ (\code{i} est entier, \code{i < 5} donne
    \code{i <= 4}), la sortie $[5,+\infty]$. Mais l'en-tête garde
    $[0,+\infty]$, et c'est ce que lit l'assertion du premier test.
  \item Avec le seuil 5, la borne saute à 5 au lieu de $+\infty$ ; la sortie
    de boucle reçoit alors exactement $[5,5]$, la valeur concrète. C'est ce
    que vérifie le second test, et ce qui permet à \code{LocalLoop}
    d'écrire $[5,5]$ dans \code{total}.
  \item Au pas 12, le corps renvoie $[1,5]$, contenu dans l'entrée de
    l'en-tête : \code{widen_to} ne touche pas une borne qui ne croît pas, la
    nouvelle entrée est égale à l'ancienne, rien n'est enfilé. La passe se
    termine dès que la file est vide.
\end{itemize}

\begin{decision}[ADR-014]
  L'ADR-014 prévoyait deux mécanismes : le widening à seuils (partie 1) et un
  \emph{narrowing} classique, c'est-à-dire une phase descendante après
  convergence (partie 2). La révision du 28 juin 2026 abandonne la partie 2.
  Constat : sur le compteur gardé, le raffinement de garde pendant la phase
  ascendante suffit ; sur la boucle locale ci-dessus, ce qui manquait était
  le widening à seuils \emph{sur l'arc arrière interne} d'\code{analyze_cfg},
  qui donne $[0,5]$ en tête et $[5,5]$ en sortie sans phase descendante. Les
  deux techniques ne récupèrent que des bornes littérales, et le widening à
  seuils les récupère déjà en montant ; une descente ne pourrait de toute
  façon pas faire redescendre le store d'état, alimenté par des joins
  monotones. Alternatives refusées : la phase descendante bornée
  (\code{narrow_iterations}), la séparation \code{WidenOutcome} dans les
  résultats et la suppression des écritures pures (\code{effect_setter_writes})
  ; aucune n'est implémentée au commit \snapshotCommit. L'opérateur
  \code{narrow} reste une infrastructure différée, à reprendre seulement si
  l'on voulait récupérer des bornes symboliques.
\end{decision}

\begin{piege}
  Le mot \og narrowing\fg{} désigne deux choses dans le projet. Les méthodes
  \code{narrow_lt}, \code{narrow_truthy}, etc. et \code{narrow_env_for_branch}
  sont du \emph{raffinement de garde} : elles restreignent une variable sous
  la condition d'une branche. Le narrowing de l'ADR-014 est une phase
  descendante du point fixe ; il n'existe pas dans le code. Le texte de
  l'ADR fait lui-même la distinction.
\end{piege}

\subsection{L'algorithme complet}

L'\cref{alg:cfg-analyzer-dominance-worklist} récapitule \code{analyze_cfg}
avec tous ses détails. Il raffine
l'\cref{alg:interpretation-abstraite-worklist} sur trois points : le store
d'état partagé, les effets de bord du \code{Return}, et le compteur par
en-tête.

\begin{algorithm}[htbp]
\KwIn{CFG d'entrée $e$, environnement $\rho_0$, stores $S$ (état) et $M$ (mémos), tas $H$, seuil $k$, seuils $T$}
\KwOut{$\mathit{Out}$ (bloc $\mapsto$ environnement de sortie), $S'$}
$\mathit{In} \gets \{e \mapsto \rho_0\}$ ; $\mathit{Out} \gets \emptyset$ ; $S' \gets S$ ; $C \gets \emptyset$ (compteurs)\;
$W \gets [e]$\;
\While{$W$ non vide}{
  $b \gets$ défiler($W$)\;
  $\rho \gets \mathit{In}[b]$ ; $M_b \gets$ copie de $M$\;
  \ForEach{instruction $i$ de $b$}{$\rho \gets \mathtt{exec\_stmt}(i, \rho)$ \tcp*{lit et écrit $S'$, $H$, $M_b$}}
  \lIf{$b$ se termine par $\mathtt{Return}(x)$}{$\mathtt{exec\_expr\_effects}(x, \rho)$}
  $\mathit{Out}[b] \gets \rho$\;
  \ForEach{$(s, \rho_s) \in \mathtt{outgoing}(b, \rho)$ \tcp*{garde raffinée sur un Branch}}{
    \uIf{$s \notin \mathrm{dom}(\mathit{In})$}{$\nu \gets \rho_s$\;}
    \uElseIf{$b \to s$ est un arc \texttt{Back}}{
      $C[s] \gets C[s] + 1$\;
      \leIf{$C[s] \ge k$}{$\nu \gets \mathit{In}[s] \widen_T \rho_s$}{$\nu \gets \mathit{In}[s] \join \rho_s$}
    }
    \lElse{$\nu \gets \mathit{In}[s] \join \rho_s$}
    \lIf{$\nu \neq \mathit{In}[s]$}{$\mathit{In}[s] \gets \nu$ ; enfiler $s$ s'il n'est pas dans $W$}
  }
}
\Return{$(\mathit{Out}, S')$}
\caption{\code{analyze_cfg} (\src{src/engine/cfg_analyzer.rs}{34}{124})}
\label{alg:cfg-analyzer-dominance-worklist}
\end{algorithm}

% ===========================================================================
\section{Raffiner par les gardes}
\label{sec:cfg-analyzer-dominance-gardes}

\subsection{Le problème}

Reprenons le compteur gardé de l'ADR-014 :

\begin{lstlisting}[language=TypeScript,caption={Un compteur gardé (\file{/tmp/ch12/guarded.tsx})},label={lst:cfg-analyzer-dominance-guarded}]
import { useState, useEffect } from "react";

export function Guarded() {
  const [count, setCount] = useState(0);
  useEffect(() => {
    if (count < 10) setCount(count + 1);
  }, [count]);
  return <div>{count}</div>;
}
\end{lstlisting}

L'effet s'exécute, incrémente, se ré-exécute, jusqu'à \code{count = 10}, puis
s'arrête. Pour le prouver, l'analyse doit utiliser la condition du \code{if} :
dans la branche \code{then}, \code{count} est strictement inférieur à 10, donc
\code{setCount} écrit au plus 10. Sans cette information, la valeur écrite
suivrait la valeur lue et croîtrait sans fin. Le moteur l'utilise :

\begin{sortie}
$ reactant check --verbose --info --no-color --project plain guarded.tsx
  [verbose] Guarded: 3 iteration(s), widened: [0]
  Guarded  (2 hooks)  guarded.tsx
    info   widening-info  (line 5:2)  state `count` kept changing during analysis and was approximated to converge, so findings that depend on it may be imprecise
    ...
    verified  infinite-loop  no effect diverges into an infinite render loop
\end{sortie}
\noindent(sortie élaguée : les autres lignes \code{verified} sont omises.)

La boucle externe a dû accélérer (\code{widened: [0]}), mais le widening à
seuils s'arrête sur le seuil 10 et le raffinement borne l'écriture ;
\regle{infinite-loop} se déclare vérifiée. Le déroulé tour par tour appartient
au \cref{chap:point-fixe-composant}. Il suffit de changer la forme de la garde
pour perdre cette preuve :

\begin{lstlisting}[language=TypeScript,caption={La même garde, écrite \code{limit > count} (\file{/tmp/ch12/unguarded.tsx})},label={lst:cfg-analyzer-dominance-unguarded}]
import { useState, useEffect } from "react";

export function Unguarded({ limit }: { limit: number }) {
  const [count, setCount] = useState(0);
  useEffect(() => {
    if (limit > count) setCount(count + 1);
  }, [count, limit]);
  return <div>{count}</div>;
}
\end{lstlisting}

\begin{sortie}
  Unguarded  (2 hooks)  unguarded.tsx
    warn   infinite-loop  [hook:0]  (line 5:2)  this effect keeps pushing state `count` (its deps do not provably gate it, so the effect can re-run every render) to new values on every run. Potential infinite render loop
\end{sortie}

La condition compare deux variables. Aucun motif ne la reconnaît, aucun seuil
ne vaut \code{limit} : \code{count} est widené jusqu'à $+\infty$. Le
\Warning{} est un défaut \emph{possible} (une prop \code{limit} valant
\code{Infinity} bouclerait vraiment), mais pour toute limite finie c'est un
faux positif, toléré par l'invariant de soundness ; le corriger demanderait un
domaine relationnel.

\subsection{Ce que fait le code}

\begin{definition}{Raffinement de garde}{cfg-analyzer-dominance-raffinement}
  Le \terme{raffinement de garde} d'un environnement $\rho$ par une condition
  $c$ sur la branche $t \in \{\mathit{vrai}, \mathit{faux}\}$ est un
  environnement $\rho'$ qui ne diffère de $\rho$ qu'en une variable $x$
  testée par $c$, avec $\rho'(x) \lleq \rho(x)$ et
  $\gam(\rho'(x)) \supseteq \{ v \in \gam(\rho(x)) \mid c[x \mapsto v]
  \text{ peut s'évaluer à } t \}$. C'est \code{narrow_env_for_branch}.
\end{definition}

La fonction commence par un utilitaire qui remplace une seule variable :

\rustsnippet[label={lst:cfg-analyzer-dominance-narrow-head}]{src/engine/cfg_analyzer.rs}{201}{217}{\code{narrow_env_for_branch} : une variable à gauche, sinon rien}

Puis elle distingue la forme du membre droit. La comparaison numérique à un
littéral a été montrée au \cref{chap:treillis-simples}
(\cref{lst:treillis-simples-narrow-env}) ; voici la suite, les comparaisons à
\code{null} et \code{undefined} et les tests de vérité :

\rustsnippet[label={lst:cfg-analyzer-dominance-narrow-null}]{src/engine/cfg_analyzer.rs}{242}{293}{raffinements de nullité et de vérité}

Le \cref{tab:cfg-analyzer-dominance-motifs} résume la fonction ; il complète
le \cref{tab:treillis-simples-motifs} par la colonne \og ce qui est
touché\fg{} et par les formes non reconnues.

\begin{table}[htbp]\centering\small
\begin{tabular}{@{}p{3.6cm}p{3.1cm}p{3.1cm}p{4.2cm}@{}}
\toprule
condition & branche vraie & branche fausse & composantes de \code{StateValue} touchées \\
\midrule
\code{x < c}, \code{x <= c}, \code{x > c}, \code{x >= c} ($c$ littéral entier ou flottant) & \code{narrow_lt}, \code{_leq}, \code{_gt}, \code{_geq} & la négation : \code{_geq}, \code{_gt}, \code{_leq}, \code{_lt} & \code{num} seul \\
\code{x == c}, \code{x != c} ($c$ numérique) & \code{narrow_eq} / \code{narrow_neq} & \code{narrow_neq} / \code{narrow_eq} & \code{num} seul \\
\code{x == null}, \code{x != null} & \code{narrow_keep_nullish_only} / \code{narrow_drop_null} & \code{narrow_drop_null} / \code{narrow_keep_nullish_only} & tout (\code{keep}) ou \code{null} (\code{drop}) \\
\code{x == undefined}, \code{x != undefined} & \code{narrow_keep_nullish_only} / \code{narrow_drop_undef} & \code{narrow_drop_undef} / \code{narrow_keep_nullish_only} & tout ou \code{undef} \\
\code{x} & \code{narrow_truthy} & \code{narrow_falsy} & toutes sauf \code{other} \\
\code{!x} & \code{narrow_falsy} & \code{narrow_truthy} & toutes sauf \code{other} \\
\midrule
\code{c < x}, \code{x < y}, \code{x.f < c}, \code{x === "a"}, \code{a && b}, \code{true}, appel… & \multicolumn{2}{l}{environnement inchangé} & aucune \\
\bottomrule
\end{tabular}
\caption{Motifs de \code{narrow_env_for_branch}
(\src{src/engine/cfg_analyzer.rs}{201}{294}). \code{==} et \code{===} sont
confondus par l'IR.}
\label{tab:cfg-analyzer-dominance-motifs}
\end{table}

\subsection{Pourquoi ces raffinements sont sound}

Trois précautions rendent ce tableau sound ; toutes sont écrites dans le
code.
\begin{itemize}
  \item \emph{Le raffinement numérique ne touche que \code{num}.} Une
    comparaison JavaScript convertit ses opérandes : \code{null < 5} est
    vrai. Retirer \code{null} ou les chaînes de la branche \og vrai\fg{}
    serait donc faux ; le commentaire de
    \src{src/domains/impls/state_value.rs}{604}{605} le dit.
  \item \emph{L'IR confond \code{==} et \code{===}.} Le raffinement doit donc
    être correct pour les deux sémantiques à la fois
    (\src{src/engine/cfg_analyzer.rs}{193}{196}). La branche où
    \code{x == null} est vrai garde \code{null} \emph{et} \code{undefined}
    (vrai de \code{==}, sur-approximation de \code{===}) ; la branche où
    \code{x != null} est vrai ne retire que \code{null} (vrai de
    \code{!==} ; pour \code{!=}, qui retirerait aussi \code{undefined},
    c'est une sur-approximation). \code{narrow_keep_nullish_only} réactive en
    outre \code{null} et \code{undefined} si la composante \code{other} est
    présente, car elle peut les cacher
    (\src{src/domains/impls/state_value.rs}{551}{558}).
  \item \emph{Un raffinement entier n'est appliqué qu'aux entiers.} Le
    passage de \code{i < 5} à \code{i <= 4} suppose que \code{i} est
    entier, ce que garantit le bit \code{is_int} ; sur un flottant la borne
    reste 5 (\cref{thm:treillis-simples-narrow}).
\end{itemize}
Enfin, tout motif non reconnu rend l'environnement inchangé, ce qui est
toujours sound : ne rien apprendre ne perd aucune valeur.

\subsection{Une garde réfutée n'élimine pas sa branche}

Que se passe-t-il quand le raffinement rend $\Bot$ ? Le composant suivant est
sain : \code{k} vaut 3, la condition \code{k > 10} est toujours fausse,
l'effet n'écrit jamais.

\begin{lstlisting}[language=TypeScript,caption={Une branche morte (\file{/tmp/ch12/dead.tsx})},label={lst:cfg-analyzer-dominance-dead}]
import { useState, useEffect } from "react";

export function Dead() {
  const [n, setN] = useState(0);
  const k = 3;
  useEffect(() => {
    if (k > 10) {
      setN(n + 1);
    }
  });
  return <div>{n}</div>;
}
\end{lstlisting}

\begin{sortie}
$ reactant check --verbose --no-color --project plain dead.tsx
  [verbose] Dead: 5 iteration(s), widened: [0]
  Dead  (2 hooks)  dead.tsx
    warn   infinite-loop  [hook:0]  (line 6:2)  this effect keeps pushing state `n` to new values on every run. Potential infinite render loop
\end{sortie}

Sur la branche vraie, \code{narrow_gt(10)} appliqué à $k = [3,3]$ donne bien
un intervalle vide. Mais \code{outgoing} rend \emph{toujours} les deux
successeurs d'un \code{Branch} (lignes 135 à 140), le successeur est enfilé
dès que son entrée change (lignes 114 à 118), et ni \code{analyze_cfg} ni le
transfert ne testent \og une variable vaut $\Bot$, donc le bloc est
inatteignable\fg. Le bloc est exécuté ; seule \code{k} y vaut $\Bot$ ;
\code{n} garde sa valeur et \code{setN(n + 1)} écrit une vraie valeur. Le
\Warning{} est un faux positif : imprécis, sound.

\begin{piege}
  Il serait tentant de \og corriger\fg{} en ne propageant pas vers un
  successeur dont l'environnement raffiné contient une variable $\Bot$. Ce
  correctif central n'est pas fait, et sa soundness n'est pas évidente : il
  faut que $\Bot$ signifie \og aucune valeur concrète\fg{} pour \emph{tout}
  le produit \code{StateValue}, y compris quand la composante \code{other}
  ou le raffinement de vérité (qui laisse \code{other} intact) interviennent.
  Aucune issue ne suit ce point au commit \snapshotCommit{} (recherche faite
  par le dossier 06, \code{gh issue list} sans résultat sur ce sujet).
\end{piege}

\subsection{Les autres utilisateurs du raffinement}

\code{narrow_env_for_branch} est \code{pub(crate)} parce que deux preuves du
moteur la réutilisent hors de toute passe : les preuves de convergence des
gardes du graphe de churn (\srcline{src/engine/guards.rs}{235} et
\srcline{src/engine/guards.rs}{771}) appliquent les gardes d'un site
d'écriture à l'environnement de sortie du rendu. Ce partage est la règle 3
du projet appliquée : une seule définition de ce qu'une garde apprend, au
niveau central (\cref{chap:churn}).

% ===========================================================================
\section{L'état dans une passe : un store insensible au flot}
\label{sec:cfg-analyzer-dominance-etat}

Les environnements sont suivis bloc par bloc ; l'état React, le tas et les
mémos ne le sont pas. Il faut comprendre pourquoi ce mélange est correct.

\paragraph{Le store d'état.} \code{state_out} est créé par
\code{state.clone()} (ligne 49), prêté à chaque bloc, et chaque setter y fait
une mise à jour faible (un join, \cref{chap:transfert-stores}). Une lecture
\code{StateVal(l)} dans un bloc voit donc les écritures de \emph{tous} les
blocs déjà visités pendant la passe, sur n'importe quel chemin. Le store est
\emph{insensible au flot} à l'intérieur d'un CFG.

\begin{react}[l'état ne change pas au milieu d'un rendu]
  Pendant un rendu, \code{count} est une constante : un \code{setCount}
  n'agit que sur la file de mises à jour, et la nouvelle valeur n'est vue
  qu'au rendu suivant. Dans un effet ou un handler, la variable \code{count}
  est de même la valeur capturée par la closure. Seul un \emph{functional
  updater} (\code{setCount(c => c + 1)}) lit l'état courant de la file. Un
  store unique, sans position, est donc proche de la sémantique de React :
  ce qui compte est l'ensemble des valeurs écrites, pas l'endroit du corps où
  on les écrit.
\end{react}

\begin{invariant}{Une passe ne fait que monter l'état}{cfg-analyzer-dominance-state-monotone}
  À tout moment d'une passe, \code{state_out} $\lgeq$ \code{state}, et
  \code{state_out} ne fait que croître. Le store rendu contient donc l'état
  d'entrée et toutes les valeurs écrites pendant la passe.
\end{invariant}

L'ordre de visite a une conséquence : un bloc visité \emph{avant} un bloc qui
écrit ne voit pas cette écriture pendant la passe. Ce n'est pas un défaut de
soundness, grâce à la boucle externe : le tour suivant repart d'un état qui
inclut l'écriture, et le point fixe n'est atteint que lorsqu'aucune passe
n'ajoute plus rien. À ce moment, \code{state_out} reste égal à \code{state}
pendant toute la passe (il part de \code{state}, ne fait que croître, et finit
$\lleq$ \code{state}), et chaque bloc a vu l'état final.

\paragraph{Le tas.} Le tas est emprunté mutablement et n'est jamais recopié :
il est partagé par toutes les passes d'un tour, et d'un tour à l'autre. C'est
voulu : une closure allouée pendant le rendu doit être visible quand un effet
l'appelle par une variable (le motif B5, commentaire
\src{src/engine/cfg_analyzer.rs}{25}{28}). Le test
\code{heap_persists_across_render_and_effect_passes} du point fixe l'épingle
(\src{src/engine/fixpoint.rs}{2157}{2232}).

\paragraph{Les mémos.} À l'inverse, le store de mémos est recopié par bloc et
les écritures locales sont jetées : une passe lit les mémos du tour, elle ne
les modifie pas. Leur recalcul est l'affaire de la boucle externe
(\cref{chap:point-fixe-composant}).

% ===========================================================================
\section{Relire une passe terminée}
\label{sec:cfg-analyzer-dominance-relire}

\subsection{Ce que garantit une passe}

Quand la file est vide, que sait-on ?

\begin{proposition}{Une passe calcule un post-point fixe des équations de blocs}{cfg-analyzer-dominance-post-point-fixe}
  Supposons que \code{state_out} ne change pas pendant la passe (c'est le cas
  au point fixe de la boucle externe). À la fin d'\code{analyze_cfg}, pour
  tout bloc visité $b$ : $\mathit{Out}[b] = f_b(\mathit{In}[b])$, où $f_b$
  est le transfert du bloc ; et pour tout successeur $(s, \rho_s)$ rendu par
  \code{outgoing} depuis $\mathit{Out}[b]$ : $\rho_s \lleq \mathit{In}[s]$.
  Si les transferts et les raffinements sont corrects, $\gam(\mathit{Out}[b])$
  contient tous les environnements concrets à la sortie de $b$.
\end{proposition}

\noindent\emph{Preuve.} Toute modification de $\mathit{In}[b]$ enfile $b$ ;
la dernière visite de $b$ a donc lieu avec l'entrée finale, et
$\mathit{Out}[b]$ est le transfert de cette entrée. Lors de cette visite,
chaque successeur a reçu $\nu \lgeq \rho_s$ (première arrivée, join ou
widening, tous au-dessus de $\rho_s$), et les entrées ne font ensuite que
croître. Le système d'équations $\mathit{In}[s] \lgeq
\bigsqcup_{p} \rho_{p \to s}$, $\mathit{Out}[b] = f_b(\mathit{In}[b])$ est
donc satisfait en post-point fixe, et le \cref{thm:post-point-fixe} donne la
sur-approximation. \qed

L'hypothèse sur \code{state_out} n'est pas décorative : pendant un tour qui
n'est pas le dernier, un bloc peut avoir été évalué contre un état plus petit
que l'état final de la passe (\cref{sec:cfg-analyzer-dominance-etat}). Seule
la boucle externe rend le résultat final sound.

\subsection{\code{entry_env_of} : reconstruire l'entrée d'un bloc}

Les règles et les relations ont parfois besoin de l'environnement
\emph{au milieu} d'un bloc, par exemple à l'instruction d'un appel de setter.
La table $\mathit{In}$ n'est pas publiée ; \code{entry_env_of} la reconstitue
à partir des sorties :

\rustsnippet[label={lst:cfg-analyzer-dominance-entry-env-of}]{src/engine/cfg_analyzer.rs}{152}{185}{relire l'environnement d'entrée d'un bloc}

Pour le bloc d'entrée, c'est l'environnement initial. Pour un autre bloc, on
rejoue \code{outgoing} depuis la sortie de chaque prédécesseur (le raffinement
de garde est donc inclus) et l'on joint. Un prédécesseur sans sortie n'a pas
été atteint et ne contribue rien. Le commentaire (lignes 156 à 159) donne
l'argument de soundness : chaque sortie sur-approxime les sorties concrètes de
son bloc, donc leur join sur-approxime l'entrée concrète. Ce join n'est jamais
widené ; il est en général \emph{plus précis} que l'entrée sur laquelle la
passe a convergé (qui peut avoir sauté à un seuil), jamais moins sûr.

L'unique consommateur est la colonne \code{written} de la relation des
écrivains, dans \src{src/engine/written.rs}{255}{274} : l'argument d'une écriture
y est évalué dans l'environnement d'entrée de son bloc, puis on rejoue les
instructions qui précèdent l'appel (\cref{chap:relations}).

\begin{exemple}{Relire l'en-tête de la boucle}{cfg-analyzer-dominance-relire}
  Sur \code{counting_loop_cfg}, les prédécesseurs de l'en-tête B1 sont B0,
  qui lui passe $[0,0]$, et B2, dont la dernière sortie vaut $[1,5]$ dans
  les deux variantes (\cref{tab:cfg-analyzer-dominance-iterations}).
  \code{entry_env_of} rend donc $[0,0] \join [1,5] = [0,5]$ pour B1, avec ou
  sans seuil. Avec $T = \{5\}$, c'est l'entrée finale de la passe. Sans
  seuil, l'entrée finale était $[0,+\infty]$ : la relecture est strictement
  plus précise que l'entrée widenée, et reste sound (\code{i} vaut
  concrètement 0 à 5 en tête). Pour B2, elle rejoue la branche vraie depuis
  $\mathit{Out}[B1]$ et rend $[0,4]$ dans les deux cas : \code{join} est le
  \code{hull} des bornes (\src{src/domains/impls/interval.rs}{71}{78}), et
  \code{narrow_lt} plafonne un entier à $v - 1$
  (\src{src/domains/impls/interval.rs}{262}{268}), ce qui rend ce calcul
  vérifiable directement sur ces deux définitions, sans exécuter
  \code{entry_env_of} elle-même.
\end{exemple}

% ===========================================================================
\section{Dominance : ce qui s'exécute sur tous les chemins}
\label{sec:cfg-analyzer-dominance-dominance}

\subsection{Le problème}

Tout ce qui précède calcule des \emph{valeurs}. Certaines règles posent une
question de \emph{structure} : ce bloc s'exécute-t-il à chaque passage dans le
corps ? Trois composants la rendent concrète :

\begin{lstlisting}[language=TypeScript,caption={Trois questions de dominance (\file{/tmp/ch12/early.tsx})},label={lst:cfg-analyzer-dominance-early}]
import { useState } from "react";

export function Panel({ ready }: { ready: boolean }) {
  if (!ready) {
    return null;
  }
  const [open, setOpen] = useState(false);
  return <div onClick={() => setOpen(!open)}>{String(open)}</div>;
}

export function Tabs({ tab }: { tab: number }) {
  const [seen, setSeen] = useState(0);
  if (tab > 0) {
    setSeen(tab);
  }
  return <div>{seen}</div>;
}

export function Always({ tab }: { tab: number }) {
  const [seen, setSeen] = useState(0);
  setSeen(tab);
  return <div>{seen}</div>;
}
\end{lstlisting}

\begin{sortie}
$ reactant check --no-color --project plain early.tsx
  Always  (1 hooks)  early.tsx
    error  setter-in-render  [hook:0]  (line 21:2)  setter `setSeen` called directly in the render body, move this call into a useEffect or an event handler
  Panel  (2 hooks)  early.tsx
    error  conditional-hook  [hook:0]  (line 7:8)  this hook is called conditionally (not on every render path)
  Tabs  (1 hooks)  early.tsx
    warn   setter-in-render  [hook:0]  (line 14:4)  setter `setSeen` called directly in the render body, move this call into a useEffect or an event handler
\end{sortie}

Dans \code{Panel}, le bloc du \code{useState} ne domine pas la sortie du
\code{return null} : le hook est sauté sur un chemin, \Error{}. Dans
\code{Tabs} et \code{Always}, la même écriture pendant le rendu reçoit deux
niveaux différents : dans \code{Always}, le bloc de l'appel domine toutes les
sorties, l'appel a lieu à chaque rendu et le défaut est certain (\Error) ;
dans \code{Tabs}, il ne les domine pas, le défaut n'est que possible
(\Warning). La règle le dit en commentaire
(\src{src/rules/impls/setter_in_render.rs}{104}{107}). La dominance est donc
l'un des outils qui \emph{certifient} un niveau \Error{}
(\cref{chap:api-regles}).

\subsection{\code{compute_dominators}}

La \cref{def:dominance} caractérise $\mathrm{Dom}(b)$ comme la plus grande
solution de $\mathrm{Dom}(e) = \{e\}$, $\mathrm{Dom}(b) = \{b\} \cup
\bigcap_{p \in \pred(b)} \mathrm{Dom}(p)$. Le code est l'algorithme itératif
classique :

\rustsnippet[label={lst:cfg-analyzer-dominance-dominators}]{src/engine/dominance.rs}{6}{58}{le plus grand point fixe de la dominance}

\begin{itemize}
  \item L'entrée est initialisée à $\{e\}$, tout autre bloc à l'ensemble de
    \emph{tous} les blocs (lignes 10 à 20) : on part du haut du treillis
    $(\powerset(\mathit{Blocs}), \supseteq)$ et l'on descend.
  \item Les blocs sont parcourus en ordre postfixe inverse (RPO), jusqu'à ce
    qu'un tour complet ne change rien (\code{changed}).
  \item Le nouvel ensemble d'un bloc est l'intersection des ensembles de ses
    prédécesseurs, plus lui-même (lignes 36 à 48). Un bloc sans
    prédécesseur autre que l'entrée est sauté (ligne 32).
  \item Les prédécesseurs sont lus dans \code{cfg.edges} par
    \code{cfg.predecessors} (\src{src/ir/cfg.rs}{173}{179}), un parcours
    linéaire des arêtes.
\end{itemize}

L'ordre RPO vient d'un parcours en profondeur récursif depuis l'entrée :

\rustsnippet[label={lst:cfg-analyzer-dominance-rpo}]{src/engine/dominance.rs}{113}{135}{ordre postfixe inverse}

Un bloc inatteignable n'est pas dans le RPO. Il n'est jamais recalculé et garde
\og tous les blocs\fg{} : \code{dominates(a, b)} est vrai pour tout $a$. C'est
neutre pour un fait must (le bloc ne s'exécute jamais), mais piégeux pour qui
énumère les blocs sans test d'atteignabilité (plus bas). Quand un prédécesseur
inatteignable participe à une intersection, son ensemble complet est l'élément
neutre de l'intersection : il ne retire rien.

\begin{figure}[htbp]\centering
\begin{tikzpicture}[node distance=9mm and 12mm]
  \node[cfgbloc] (b0) {B0 (entrée)\\branch};
  \node[cfgbloc,below left=of b0] (b1) {B1\\jump B3};
  \node[cfgbloc,below right=of b0] (b2) {B2\\jump B3};
  \node[cfgbloc,below=27mm of b0] (b3) {B3 (en-tête)\\branch};
  \node[cfgbloc,below left=of b3] (b4) {B4 (corps)\\jump B3};
  \node[cfgbloc,below right=of b3] (b5) {B5\\return};
  \node[cfgbloc,right=22mm of b2,draw=ra-amber] (b6) {B6 (orphelin)\\return};
  \draw[fleche] (b0) -- (b1);
  \draw[fleche] (b0) -- (b2);
  \draw[fleche] (b1) -- (b3);
  \draw[fleche] (b2) -- (b3);
  \draw[fleche] (b3) -- (b4);
  \draw[fleche] (b3) -- (b5);
  \draw[flechemay] (b4.west) to[bend left=45] node[etiquette,left]{Back} (b3.west);
  \node[etiquette,above=1mm of b0] {$\{0\}$};
  \node[etiquette,left=1mm of b1] {$\{0,1\}$};
  \node[etiquette,right=1mm of b2] {$\{0,2\}$};
  \node[etiquette,right=2mm of b3] {$\{0,3\}$};
  \node[etiquette,below=1mm of b4] {$\{0,3,4\}$};
  \node[etiquette,below=1mm of b5] {$\{0,3,5\}$};
  \node[etiquette,below=1mm of b6,align=center] {$\{0,\dots,6\}$\\(jamais recalculé)};
\end{tikzpicture}
\caption{Un losange suivi d'une boucle, et un bloc orphelin. Sous chaque bloc,
son ensemble de dominateurs, tel que \code{compute_dominators} le rend (relevé
à la sonde). RPO : $[0, 2, 1, 3, 5, 4]$.}
\label{fig:cfg-analyzer-dominance-dom}
\end{figure}

\begin{exemple}{Dominateurs d'un losange suivi d'une boucle}{cfg-analyzer-dominance-dom}
  Sur la \cref{fig:cfg-analyzer-dominance-dom} (un rendu schématique du type
  \code{if (fast) {...} else {...} while (delay > 10) {...} return ...}, plus un
  \code{return} orphelin comme en laisse un \code{if}/\code{else} dont les
  deux branches retournent), la sonde donne le RPO $[0,2,1,3,5,4]$. Premier
  tour : $\mathrm{Dom}(2) = \{0,2\}$, $\mathrm{Dom}(1) = \{0,1\}$ ;
  $\mathrm{Dom}(3) = \{3\} \cup (\{0,1\} \cap \{0,2\} \cap \mathit{Tous}) =
  \{0,3\}$ (B4, visité plus tard, vaut encore \og tous\fg{} et ne retire
  rien) ; $\mathrm{Dom}(5) = \{0,3,5\}$ ; $\mathrm{Dom}(4) = \{0,3,4\}$.
  Second tour : $\mathrm{Dom}(3) = \{3\} \cup (\{0,1\} \cap \{0,2\} \cap
  \{0,3,4\}) = \{0,3\}$, inchangé ; rien ne bouge, l'algorithme s'arrête.
  B6, inatteignable, garde $\{0,\dots,6\}$.
\end{exemple}

\subsection{\code{DominatorTree} et \code{dominates}}

Recalculer la dominance à chaque question coûte un point fixe par question.
\code{DominatorTree} mémorise le résultat :

\rustsnippet[label={lst:cfg-analyzer-dominance-domtree}]{src/engine/dominance.rs}{94}{111}{la relation précalculée}

Malgré son nom, la structure ne contient pas un arbre de dominateurs
immédiats mais la table complète des ensembles. La question
\code{dominates(a, b)} est un test d'appartenance ; un bloc $b$ absent de la
table (identifiant inconnu) répond faux. La fonction libre
\code{dominates(cfg, a, b)} (\src{src/engine/dominance.rs}{60}{67}) reconstruit
tout à chaque appel, et sa documentation met en garde : \og otherwise this is
O(queries × fixpoint)\fg.

\subsection{Les sorties atteignables : \code{ExitDominance}}

\og S'exécute à chaque rendu\fg{} signifie \og domine toutes les sorties\fg.
Encore faut-il savoir ce qu'est une sortie. La couche règles en a un
propriétaire unique :

\rustsnippet[label={lst:cfg-analyzer-dominance-exitdom}]{src/rules/api/query.rs}{647}{671}{les sorties sont les \code{Return} atteignables}

Deux conventions s'y lisent. Une sortie est un bloc terminé par
\code{Return} ; un \code{throw} (terminateur \code{Unreachable}) n'en est pas
une. Et un \code{Return} que rien n'atteint n'est pas une sortie : le filtre
passe par \code{CFG::reachable_blocks} (\src{src/ir/cfg.rs}{153}{163}). Les
méthodes \code{certify} et \code{may_be_skipped} qui suivent ont été
présentées au \cref{chap:interpretation-abstraite}
(\cref{lst:interpretation-abstraite-exitdom}).

\begin{decision}[ADR-025]
  Avant l'ADR-025, la fin d'un corps sans \code{return} était scellée
  \code{Unreachable}, comme un \code{throw}. Une greffe (\cref{chap:splice-inlining})
  ne raccordait donc pas un appelé sans \code{return} à la suite de
  l'appelant : le \code{Return} de l'appelant n'avait plus de prédécesseur,
  \code{block_states} n'en avait pas de sortie, et \code{exit_env} joignait
  \emph{moins de chemins que le programme n'en a} (208 composants touchés sur
  le corpus, dans la direction interdite). Décision : (1) une fin de corps
  est scellée \code{Return(Lit(Unit))} ; (2) \code{throw} garde
  \code{Unreachable} et la greffe continue de l'ignorer ; (3) un
  \code{Return} inatteignable n'est pas une sortie. Alternative refusée pour
  (2) : relier chaque \code{Unreachable} de l'appelé à la jonction. Elle
  invente un chemin qui contourne les hooks de l'appelé et produisait un
  \regle{conditional-hook} \Error{} sur l'idiome conforme
  \code{if (ctx === undefined) throw ...; useState(...)}. Le point (3) est la
  conséquence du point (1) : un \code{if}/\code{else} dont les deux
  branches retournent laisse désormais une queue \code{Return} orpheline, et
  la compter rendrait \og conditionnel\fg{} tout hook placé avant.
\end{decision}

\begin{piege}
  B6, dans la \cref{fig:cfg-analyzer-dominance-dom}, est un \code{Return}
  que rien n'atteint. La table le déclare dominé par \emph{tous} les blocs,
  B4 compris, alors qu'aucun chemin ne passe par B4 avant B6 : ce n'est pas
  un fait sur le programme, c'est un artefact de l'initialisation. Toute
  réponse de \code{dominates(_, 6)} est donc vide de sens, et un
  consommateur qui énumère \og tous les \code{Return}\fg{} sans
  \code{reachable_blocks} fait dépendre son verdict de cet artefact.
  \code{ExitDominance} ne pose jamais la question sur une sortie
  inatteignable ; c'est la seule façon correcte d'interroger la dominance
  des sorties.

  Ce constat rend d'ailleurs difficile à rejouer, avec le code actuel, le
  mécanisme que décrit le point (3) de l'\adr{25} : une queue \code{Return}
  orpheline comptée parmi les sorties ferait \emph{échouer} tout hook placé
  avant la branche à \og dominer toutes les sorties\fg. Or, avec
  \code{compute_dominators} tel qu'il est écrit au commit \snapshotCommit,
  c'est l'inverse qui se lit sur B6 : un bloc inatteignable est dominé par
  \emph{tous} les blocs, donc \code{dominates(a, b)} vaut vrai pour tout $a$
  quand $b$ est l'orphelin, et compter l'orphelin parmi les sorties ferait
  dominer \emph{trivialement} cette sortie fantôme par tout hook, sans faire
  échouer \code{certify}. L'\adr{25} ne précise pas si
  \code{compute_dominators} avait, à sa date (29 juillet 2026), la même
  convention d'initialisation qu'au commit \snapshotCommit{} ; le mécanisme
  exact qui aurait produit le \regle{conditional-hook} \Error{} qu'elle
  décrit n'est pas reconstituable depuis le seul code présent. Le point ne
  change rien à la décision retenue : filtrer les sorties par
  \code{reachable_blocks} est correct indépendamment de ce mécanisme.
\end{piege}

\subsection{\code{on_all_paths} : un ensemble sur tous les chemins}

La seconde question porte sur un \emph{ensemble} de blocs : tout chemin de
l'entrée à une sortie passe-t-il par l'un d'eux ? C'est ce que demande
\regle{infinite-loop} quand elle veut certifier qu'un effet écrit une valeur
fraîche à chaque exécution, quelle que soit la branche prise
(\srcline{src/rules/impls/infinite_loop.rs}{454}), ou le graphe de churn pour
ses arêtes \code{Must} (\srcline{src/engine/churn.rs}{538}).

\rustsnippet[label={lst:cfg-analyzer-dominance-on-all-paths}]{src/engine/dominance.rs}{69}{92}{l'ensemble post-domine-t-il l'entrée ?}

L'idée est de chercher un contre-exemple : un chemin qui part de l'entrée,
évite l'ensemble et atteint une feuille (un bloc sans successeur). S'il en
existe un, la réponse est non ; sinon oui. Plusieurs détails comptent :
\begin{itemize}
  \item \emph{Toute feuille compte.} Un \code{Return} comme un
    \code{Unreachable} (\code{throw}) : \code{cfg.successors} est vide dans
    les deux cas. Un chemin d'exception qui évite l'ensemble suffit à dire
    non. C'est le choix prudent pour un fait must.
  \item \emph{Le parcours est un DFS.} Le commentaire dit \og BFS\fg, mais
    \code{Vec::pop} dépile le dernier élément. Le résultat est le même,
    seul l'ordre de visite change.
  \item \emph{Un cycle qui évite l'ensemble ne dit pas non.} Un chemin infini
    n'atteint aucune feuille. En particulier, sur un CFG qui ne contient
    aucune feuille atteignable (une boucle \code{while (true)} sans
    \code{break}), la fonction rend \emph{vrai pour tout ensemble} ; la
    sonde l'a vérifié avec un ensemble réduit à un bloc inexistant.
\end{itemize}

\begin{remarque}
  Cette vacuité contraste avec \code{ExitDominance}, dont la documentation
  dit qu'\og a CFG with no reachable exit proves nothing in either
  direction\fg{} et qui rend \code{None}. Les deux quantificateurs n'ont pas
  la même convention sur un corps qui ne termine jamais, ni sur
  \code{throw} (une feuille pour \code{on_all_paths}, pas une sortie pour
  \code{ExitDominance}). Chacune est défendable pour son usage ; il faut le
  savoir avant de remplacer l'une par l'autre.

  \regle{infinite-loop} appelle \code{must_on_all_paths} sur le CFG d'un
  effet réel, jamais sur un CFG construit à la main : la question est donc
  de savoir si un tel corps peut n'avoir aucune feuille atteignable. La
  réponse est non. \code{lower_while}
  (\src{src/lowering/cfg_builder.rs}{603}{632}) ajoute
  \emph{inconditionnellement} l'arête \code{IfFalse} vers le bloc de sortie
  de la boucle, que la condition soit ou non un littéral \code{true} ; la
  \cref{sec:cfg-analyzer-dominance-gardes} a montré que le raffinement de
  garde n'élague jamais une arête. Un \code{while (true)} lowé garde donc
  toujours une sortie structurellement atteignable, et le corps entier finit
  par un \code{Return} scellé par l'\adr{25}. La vacuité de
  \code{on_all_paths} n'a donc pas d'incidence connue sur
  \regle{infinite-loop} au commit \snapshotCommit{} : elle reste une
  possibilité théorique du treillis, pas un cas observé, et un CFG où elle
  se manifesterait devrait être construit à la main comme celui de la sonde.
\end{remarque}

Sur la \cref{fig:cfg-analyzer-dominance-dom}, la sonde donne :
$\{1\}$ non (le chemin par B2 l'évite), $\{1, 2\}$ oui, $\{3\}$ oui,
$\{4\}$ non (on peut sortir sans entrer dans la boucle), $\{5\}$ oui. Sur le
diamant seul (\src{src/engine/dominance.rs}{192}{258}), $\{1\}$ non,
$\{1,2\}$ oui, $\{3\}$ oui. On voit la différence avec la dominance : aucun
bloc de $\{1, 2\}$ ne domine B3, mais l'ensemble est sur tous les chemins.

\code{on_all_paths} n'est pas la seule implémentation du quantificateur
\og sur tous les chemins\fg. \code{must_setter_on_all_paths}
(\src{src/rules/api/query.rs}{527}{622}) porte sa propre analyse must en
avant sur des booléens, et teste tous les blocs \code{Return} et
\code{Unreachable}. La couche règles et sa façade typée
\code{must_on_all_paths} (\src{src/rules/api/query.rs}{624}{640}) sont
présentées au \cref{chap:api-regles}. L'ADR-042 (§6) a déplacé
\code{on_all_paths} des règles vers \file{src/engine/dominance.rs} : c'est une
propriété du CFG, et le graphe de churn, qui est un produit du moteur, en a
besoin.

% ===========================================================================
\section{Sonder après convergence : \code{eval_in_stores} et \code{ConvergedEval}}
\label{sec:cfg-analyzer-dominance-sondes}

\subsection{Le problème}

Après le point fixe, les relations du moteur et les règles ont besoin de
\emph{valeurs} : que vaut la dep \code{options} de cet effet ? l'argument de
ce setter est-il un objet frais ? Elles ne relancent pas \code{analyze_cfg} :
elles évaluent une expression dans un environnement donné (souvent
\code{exit_env}), contre les stores du composant.

\begin{definition}{Sonde post-convergence}{cfg-analyzer-dominance-sonde}
  Une \terme{sonde} est l'évaluation d'une expression de l'IR, dans un
  environnement donné, contre une \emph{copie} des stores d'un composant,
  après son point fixe. Elle ne modifie ni les stores ni le tas du résultat.
\end{definition}

\subsection{Le cœur : \code{eval_in_stores}}

\rustsnippet[label={lst:cfg-analyzer-dominance-eval-in-stores}]{src/engine/eval.rs}{31}{46}{le cœur mécanique de toute sonde}

La fonction clone l'état et les mémos, construit un \code{AnalysisCtx} nul
(\code{NullCtx}, pas d'inter-composants) et appelle \code{eval_expr}. Le tas,
lui, est emprunté mutablement : l'appelant passe un tas jetable, parce que
l'évaluation \emph{écrit} dans le tas (un \code{ObjectLit} sondé y crée une
entrée). Le choix des stores est laissé à l'appelant, et la documentation
(\src{src/engine/eval.rs}{19}{30}) insiste : les stores convergés et un tas
convergé, ou des stores vides pour une évaluation \og au montage\fg, ne sont
pas interchangeables.

\subsection{La couche commune : \code{ConvergedEval}}

\rustsnippet[label={lst:cfg-analyzer-dominance-converged}]{src/engine/eval.rs}{61}{80}{le cas commun : stores et tas convergés}

\rustsnippet[label={lst:cfg-analyzer-dominance-eval-struct}]{src/engine/eval.rs}{89}{105}{un évaluateur réutilisable}

\code{result.evaluator()} rend un \code{Eval} qui possède \emph{une} copie du
tas convergé ; \code{at} évalue une expression contre l'état, les mémos et ce
tas. \code{eval_in} est la version à un coup. Réutiliser le même
\code{Eval} pour plusieurs sondes (toutes les deps d'un effet, par exemple
\srcline{src/rules/impls/missing_deps.rs}{155}) évite un clone du tas par
expression ; c'est sûr parce qu'un \code{ExprId} nomme un seul site
d'allocation (\cref{def:site}, \issue{134}) : sonder deux fois le même site
réécrit sa propre entrée, deux sites ne se télescopent pas.

\begin{piege}
  L'histoire de cette couche est celle d'un faux négatif (\issue{135}).
  Le tas était un argument de l'appelant, au motif qu'un tas vide et un tas
  convergé étaient deux choix légitimes. Quatre sites sur six avaient choisi
  le tas vide, à tort : une lecture de membre (\code{obj.f},
  \code{form.onSubmit}) ne se résout que par le tas, un tas vide répond
  $\Top$, et $\Top$ est le côté \emph{silencieux} de presque tous les
  prédicats des règles (\og je ne sais pas si c'est frais, je ne dis
  rien\fg). Le correctif a été central et non site par site : figer le cas
  commun dans \code{ConvergedEval} (\src{src/engine/eval.rs}{48}{60}). Un
  site qui veut réellement des stores vides appelle \code{eval_in_stores}
  directement.
\end{piege}

\code{eval_in_stores} vivait dans \file{rules/helpers} ; l'ADR-042 (§6) l'a
descendue dans le moteur, parce que les relations dérivées après convergence
(fraîcheur des écrivains, déclencheurs d'effets, graphe de churn) en ont
besoin et ne peuvent dépendre de la couche règles. Le point fixe l'utilise
lui-même pour \relation{effect_triggers} (\src{src/engine/fixpoint.rs}{647}{667}).

% ===========================================================================
\section{Terminaison et complexité}
\label{sec:cfg-analyzer-dominance-terminaison}

\subsection{Pourquoi une passe termine}

\begin{invariant}{Tout cycle porte un arc arrière}{cfg-analyzer-dominance-back}
  Dans un CFG produit par le lowering, tout cycle contient au moins une arête
  \code{EdgeKind::Back}. Le lowering marque ainsi la fin du corps d'un
  \code{while}, d'un \code{for}, d'un \code{for...of}, le retour d'un
  \code{do...while} et tout \code{continue}
  (\cref{def:lowering-back} ; \srcline{src/lowering/cfg_builder.rs}{350},
  \srcline{src/lowering/cfg_builder.rs}{396}, \srcline{src/lowering/cfg_builder.rs}{629},
  \srcline{src/lowering/cfg_builder.rs}{699}, \srcline{src/lowering/cfg_builder.rs}{768}).
\end{invariant}

\begin{theoreme}{Terminaison d'une passe}{cfg-analyzer-dominance-terminaison}
  Sous l'\cref{inv:cfg-analyzer-dominance-back}, \code{analyze_cfg} termine
  sur tout CFG fini.
\end{theoreme}

\noindent\emph{Esquisse.} Chaque entrée $\mathit{In}[s]$ ne fait que croître :
elle est remplacée par un join ou un widening d'elle-même, tous deux
au-dessus d'elle (\code{widen_to} $\lgeq$ join,
\cref{thm:treillis-simples-widen-to}). Un bloc n'est enfilé que si son entrée
change strictement. Il suffit donc de montrer que chaque entrée ne change
qu'un nombre fini de fois.
\begin{itemize}
  \item Dans une valeur \code{StateValue}, seule la composante numérique est
    de hauteur infinie : booléens, nullité, stabilité, setters et
    \code{StrConst} sont de hauteur finie
    (\cref{chap:treillis-simples,chap:statevalue-stabilite}). Une variable
    qui apparaît ou disparaît d'un côté passe à $\Top$, qui est le sommet.
    Le nombre de variables est borné par le CFG.
  \item Sur l'entrée d'un en-tête, à partir de la
    \code{widen_threshold}-ième arrivée par un arc arrière, la composante
    numérique passe par \code{widen_to}, qui ne peut changer une borne
    qu'un nombre fini de fois ($T$ est fini, puis $\pm\infty$).
  \item Un bloc qui n'est pas un en-tête reçoit un join d'un nombre fini de
    prédécesseurs. Retirer les arcs arrière rend le graphe acyclique
    (invariant) : par récurrence dans l'ordre topologique de ce graphe
    acyclique, une fois que les en-têtes ne changent plus, chaque bloc ne
    reçoit plus qu'un nombre fini de valeurs différentes.
\end{itemize}
Le nombre total de changements est donc fini, et la file finit par se vider.
\qed

Il faut remarquer que les arrivées par un arc \emph{non} arrière sur un
en-tête (depuis le pré-en-tête) sont jointes, jamais widenées. C'est suffisant,
parce que le pré-en-tête est hors de la boucle : dans une boucle imbriquée,
l'entrée qu'il transmet ne croît qu'autant que la boucle extérieure, elle-même
bornée par son propre arc arrière.

\begin{piege}
  L'invariant est fragile, et sa violation ne se voit pas : un cycle dont
  aucune arête n'est marquée \code{Back} n'est jamais widené. Le cas s'est
  produit : un \code{continue} émis avec une arête \code{Unconditional}.
  Le commentaire de la fixture qui l'épingle le raconte :
  \og Marked \code{Unconditional}, the join grew \code{i} by one per pass
  and the fixpoint never terminated\fg{}
  (\src{tests/fixtures/widening.tsx}{51}{54}). Toute modification du lowering
  qui crée un cycle doit marquer son arête de retour ; toute greffe
  (\cref{chap:splice-inlining}) doit conserver la nature des arêtes.
\end{piege}

\subsection{Complexité}

Soient $N$ le nombre de blocs, $E$ le nombre d'arêtes, $h$ le nombre maximal
de changements d'une entrée (borné par la hauteur des domaines après
widening). Le \cref{tab:cfg-analyzer-dominance-complexite} résume les coûts,
tels qu'ils se lisent dans le code.

\begin{table}[htbp]\centering\small
\begin{tabular}{@{}p{3.3cm}p{5.2cm}p{5.4cm}@{}}
\toprule
fonction & coût & d'où il vient \\
\midrule
\code{analyze_cfg} & $O(h \cdot N)$ visites ; chaque visite : transfert, clone de l'entrée et des mémos, $O(E)$ par successeur & \code{is_back} parcourt \code{cfg.edges} pour chaque successeur (lignes 92 à 95) \\
\code{compute_dominators} & $O(R \cdot N \cdot (E + N))$ pour $R$ tours & \code{predecessors} en $O(E)$, intersections d'ensembles en $O(N)$ \\
\code{dominates} (méthode) & $O(1)$ attendu & test d'appartenance dans un \code{HashSet} \\
\code{dominates} (libre) & un \code{compute_dominators} par appel & à éviter en boucle \\
\code{on_all_paths} & $O(N \cdot E)$ & chaque bloc visité appelle \code{successors} en $O(E)$ \\
\code{entry_env_of} & $O(E)$ + un \code{outgoing} par prédécesseur & filtre des arêtes entrantes \\
\bottomrule
\end{tabular}
\caption{Coûts des fonctions du chapitre.}
\label{tab:cfg-analyzer-dominance-complexite}
\end{table}

Aucun de ces coûts ne pose problème sur des corps de composants, qui ont
typiquement quelques dizaines de blocs ; le seul piège documenté est la
reconstruction de la dominance par question, qui a motivé
\code{DominatorTree} (\og the fix for \code{dominates()} driving per-rule
quadratic recomputation\fg, \src{src/engine/dominance.rs}{94}{95}). Pour le
nombre de tours $R$ de \code{compute_dominators}, le parcours en RPO fait que,
sur un CFG réductible, quelques tours suffisent ; l'exemple de la
\cref{fig:cfg-analyzer-dominance-dom} en fait deux.

\subsection{Deux interpréteurs de CFG}

\begin{remarque}
  \code{analyze_cfg} n'est pas le seul interpréteur de CFG du projet.
  \code{exec_body_impl} (\file{src/domains/interp/interpreter.rs},
  \cref{chap:transfert-stores}) exécute les corps de fonctions imbriquées
  (callbacks, \code{.then}, updaters, initialiseurs) en une seule passe
  topologique, arcs arrière ignorés et sans widening : une valeur portée par
  une boucle n'y est vue qu'à sa première itération. Les deux issues ouvertes
  \issue{12} (\og Two CFG interpreters of different strength, and nothing
  says which ran\fg) et \issue{21} (faux négatif sur les valeurs portées par
  une boucle dans un callback) suivent cette asymétrie. Le handler de
  \code{onClick} passe par \code{analyze_cfg} ; la flèche passée à
  \code{setTimeout} dans un effet passe par \code{exec_body_impl}.
\end{remarque}

% ===========================================================================
\begin{resume}
\begin{itemize}
  \item \code{analyze_cfg} (\file{src/engine/cfg_analyzer.rs}) interprète un
    CFG par worklist FIFO : environnements d'entrée et de sortie par bloc,
    un seul store d'état pour toute la passe (insensible au flot, mis à jour
    faiblement), un tas partagé entre passes, des mémos figés. Elle rend les
    environnements de \emph{sortie} (malgré la documentation de
    \code{BlockEnvs}) et le store enrichi.
  \item Fusion : première arrivée, join, et \code{widen_to} à seuils sur un
    arc \code{Back} dès la \code{widen_threshold}-ième arrivée (3), compteur
    par en-tête local à la passe. Les seuils sont les littéraux numériques du
    composant. Le narrowing descendant de l'ADR-014 n'existe pas : le
    widening à seuils sur l'arc interne suffit (\code{[0,5]} en tête,
    \code{[5,5]} en sortie de la boucle \code{i < 5}).
  \item Les effets de bord d'un \code{Return} sont exécutés
    (\code{exec_expr_effects}) : une flèche concise écrit.
  \item \code{narrow_env_for_branch} raffine une variable comparée à un
    littéral, à \code{null}/\code{undefined}, ou testée pour sa vérité ;
    tout autre motif ne raffine rien. Une garde réfutée n'élimine pas sa
    branche (FP de \code{Dead}).
  \item Une passe finie est un post-point fixe des équations de blocs ;
    \code{entry_env_of} relit l'entrée d'un bloc par un join jamais widené ;
    \code{exit_env} joint les sorties des \code{Return} atteints (par
    \code{reduce} : la table vide n'est pas neutre).
  \item \code{compute_dominators} calcule les ensembles de dominateurs en
    RPO ; \code{DominatorTree} les mémorise ; \code{ExitDominance} ne compte
    que les \code{Return} atteignables (ADR-025) ; \code{on_all_paths}
    cherche un chemin qui évite un ensemble et atteint une feuille, \code{throw}
    compris.
  \item \code{eval_in_stores} est le cœur de toute sonde post-convergence ;
    \code{ConvergedEval}/\code{Eval} en figent le cas commun avec le tas
    convergé (\issue{135}).
  \item Fichiers : \file{src/engine/cfg_analyzer.rs},
    \file{src/engine/dominance.rs}, \file{src/engine/eval.rs}, et dans
    \file{src/engine/fixpoint.rs} \code{exit_env} et
    \code{collect_thresholds}.
\end{itemize}
Le \cref{chap:point-fixe-composant} enchaîne ces passes : rendu, mémos, effets,
handlers, test de convergence et widening de l'état, puis les passes d'après
convergence.
\end{resume}

% ===========================================================================
\section*{Exercices}
\addcontentsline{toc}{section}{Exercices}
\label{sec:cfg-analyzer-dominance-exercices}

\begin{exercice}{Dominateurs d'un losange}{cfg-analyzer-dominance-losange}
  Soit le losange du test \code{diamond_cfg}
  (\src{src/engine/dominance.rs}{192}{258}) : B0 se branche vers B1 et B2,
  qui sautent tous deux vers B3, terminé par \code{return}.
  \begin{enumerate}
    \item Donner un RPO possible, puis dérouler \code{compute_dominators}
      tour par tour.
    \item Que rendent \code{dominates(1, 3)}, \code{on_all_paths} sur
      $\{1\}$ et \code{on_all_paths} sur $\{1, 2\}$ ?
    \item On ajoute un bloc B4 terminé par \code{return}, sans arête
      entrante. Que devient $\mathrm{Dom}(4)$ ? \code{ExitDominance}
      compte-t-il B4 parmi les sorties ?
  \end{enumerate}
\end{exercice}
\begin{solution}
  (1) Le DFS visite B1 puis B3, puis B2 : post-ordre $[3,1,2,0]$, RPO
  $[0,2,1,3]$. Premier tour : $\mathrm{Dom}(2) = \{0,2\}$, $\mathrm{Dom}(1)
  = \{0,1\}$, $\mathrm{Dom}(3) = \{3\} \cup (\{0,1\} \cap \{0,2\}) =
  \{0,3\}$. Second tour : rien ne change. (2) Faux (test
  \code{diamond_branch_does_not_dominate_join}) ; faux ; vrai (valeurs
  relevées à la sonde). (3) B4 n'est pas dans le RPO et garde tous les
  blocs, $\{0,1,2,3,4\}$ ; \code{ExitDominance} l'écarte car il n'est pas
  dans \code{reachable_blocks}.
\end{solution}

\begin{exercice}{Une boucle avec des seuils}{cfg-analyzer-dominance-seuils}
  Dérouler \code{analyze_cfg} sur l'IR de \code{let i = 0; while (i < 3) { i = i + 2; }} avec \code{widen_threshold = 3} et $T = \{0, 2, 3\}$ (les
  littéraux du programme). Quelles sont les valeurs finales de \code{i} en
  tête et en sortie de boucle ? Comparer à la valeur concrète.
\end{exercice}
\begin{solution}
  B0 installe $\mathit{In}[B1] = [0,0]$. B1 passe $[0,0]$ au corps et un
  intervalle vide à la sortie. Le corps renvoie $[2,2]$ ; premier retour,
  join : $\mathit{In}[B1] = [0,2]$. B1 passe au corps $[0,2]$ (raffiné par
  \code{i < 3}, soit \code{i <= 2}) ; le corps renvoie $[2,4]$ ; deuxième
  retour, join : $\mathit{In}[B1] = [0,4]$. B1 raffine $[0,4]$ en $[0,2]$
  pour le corps, \emph{égal} à son entrée : le corps n'est pas ré-enfilé, et
  il n'y aura pas de troisième retour. La sortie reçoit $[3,4]$. Résultat :
  $[0,4]$ en tête, $[3,4]$ en sortie, sans aucun widening (le compteur de B1
  s'arrête à 2). La valeur concrète en sortie est 4, contenue dans $[3,4]$.
  Ce déroulé n'a pas été exécuté ; il se relit directement sur \code{add}
  (\src{src/domains/impls/interval.rs}{148}{156}), \code{hull}
  (\src{src/domains/impls/interval.rs}{71}{78}) et
  \code{narrow_lt}/\code{narrow_geq}
  (\src{src/domains/impls/interval.rs}{262}{292}).
\end{solution}

\begin{exercice}{Lire un raffinement}{cfg-analyzer-dominance-raffinement}
  Pour chaque condition, dire ce que reçoivent les deux branches si
  \code{x} vaut \og nombre $[0,10]$ ou \code{null}\fg{} : \code{x < 5},
  \code{5 > x}, \code{x != null}, \code{!x}, \code{x === 0}.
\end{exercice}
\begin{solution}
  \code{x < 5} : vrai $[0,4]$ ou \code{null} (entier ; \code{null} gardé car
  la comparaison le convertit), faux $[5,10]$ ou \code{null}. \code{5 > x} :
  motif non reconnu, les deux branches reçoivent l'environnement inchangé.
  \code{x != null} : vrai $[0,10]$ (\code{null} retiré), faux :
  \code{narrow_keep_nullish_only}, \code{null} seul. \code{!x} : vrai
  \code{narrow_falsy}, $[0,0]$ ou \code{null} ; faux \code{narrow_truthy},
  \code{null} retiré et \code{narrow_neq(0)} sur $[0,10]$, qui ne rend $\Bot$
  que sur le point $[0,0]$ et laisse tout autre intervalle intact
  (\src{src/domains/impls/interval.rs}{303}{309}) : $[0,10]$. \code{x === 0} :
  \code{narrow_eq(0)}, $[0,0]$ ou \code{null} ; faux \code{narrow_neq(0)},
  $[0,10]$ ou \code{null}.
\end{solution}

\begin{exercice}{Pourquoi \code{!=} suffit}{cfg-analyzer-dominance-egalite}
  Montrer que, dans \code{analyze_cfg}, la nouvelle entrée $\nu$ d'un bloc
  vérifie toujours $\nu \lgeq \mathit{In}[s]$, et en déduire que le test
  \code{!=} (ligne 114) équivaut à $\nu \not\lleq \mathit{In}[s]$. Quelle
  hypothèse sur l'égalité de \code{AbstractEnv} ce raisonnement demande-t-il,
  et la remarque sur les intervalles vides de la
  \cref{sec:cfg-analyzer-dominance-diamant} la contredit-elle ?
\end{exercice}
\begin{solution}
  $\nu$ est $\mathit{In}[s] \join \rho$ ou $\mathit{In}[s] \widen_T \rho$,
  tous deux $\lgeq \mathit{In}[s]$ ; pour $\nu \lgeq \mathit{In}[s]$,
  $\nu \lleq \mathit{In}[s]$ équivaut à l'égalité \emph{dans l'ordre}. Le
  raisonnement demande que l'égalité structurelle coïncide avec l'égalité
  d'ordre. Les intervalles vides la contredisent : deux représentations de
  $\Bot$ sont égales pour l'ordre et différentes pour \code{==}. La
  conséquence n'est qu'un ré-enfilage superflu (le test voit un changement
  qui n'en est pas un), jamais un changement manqué : la terminaison n'est
  pas menacée, car ces représentations suivent l'entrée de l'en-tête, qui se
  stabilise.
\end{solution}

\begin{exercice}{Une garde morte}{cfg-analyzer-dominance-dead}
  Sur \cref{lst:cfg-analyzer-dominance-dead}, expliquer pourquoi le
  raffinement $k := \Bot$ n'empêche pas l'écriture. Proposer un test qui
  ferait de la branche une branche morte, dire où le placer pour qu'il
  profite à toutes les passes et à \code{entry_env_of}, et énumérer les
  composantes de \code{StateValue} pour lesquelles il faut vérifier que
  $\Bot$ signifie \og aucune valeur concrète\fg.
\end{exercice}
\begin{solution}
  Le raffinement ne porte que sur \code{k} ; \code{n} garde sa valeur et
  \code{setN(n + 1)} écrit. Un test naturel serait dans \code{outgoing} :
  ne pas rendre un successeur dont l'environnement raffiné lie une variable
  à une valeur entièrement $\Bot$, ce qui profite à \code{analyze_cfg} et à
  \code{entry_env_of}, qui partagent \code{outgoing}. Il faut vérifier toutes
  les composantes du produit (\code{num}, \code{boolean}, \code{str},
  \code{reference}, \code{null}, \code{undef}, \code{setter}, \code{other}),
  et en particulier que \code{other}, jamais retiré par les raffinements,
  soit faux. Ce n'est pas fait au commit \snapshotCommit{}.
\end{solution}
